\chapter{Hybrid Representations for Similarity and Label Noise}\label{chap:similarity}
% ---- TAKEAWAYS ----
% - The hybrid began as a fix for similarity; this chapter tests whether features chosen for the task give better neighbours than the single representations they come from.
% - A similarity hybrid uses the same search but ranks candidates by predicting each molecule's activity from its five nearest neighbours.
% - Finding noisy labels, estimating how large the noise is, and improving the held-out score are three different tasks. A method can do the first and fail the third.
% - The type of noise decides what can be fixed. Censoring has a known direction and location and is largely correctable; Gaussian noise, the standard assumption, is the hardest case.
% - Under Gaussian noise the input moved the score more than the choice of LNL method did (preliminary).
% - Better similarity cannot remove the noise on the neighbours' own labels. No one method fits every noise type and dataset, so model the noise first.
% ---- END TAKEAWAYS ----

\section{Introduction}
% ---- NOTES: Introduction ----
% Takeaway: Chapter 4 showed where similarity breaks; Chapter 5 built a representation chosen for the task. This chapter asks whether that representation repairs similarity, noise detection and learning with noisy labels.
% Takeaway: Short, about one page: three paragraphs, the research questions as a list, and one status sentence.
% Paragraph 1: what the previous two chapters leave open
% - Chapter \ref{chap:flaws} showed where the similarity assumption fails on standard representations: ECFP4 bits compared by Tanimoto mix a few task-relevant features with many that are constant or irrelevant on the dataset.
% - Methods that trust a molecule's neighbours inherit that error, including finding noisy labels by their neighbours.
% - Chapter \ref{chap:method} built the adaptive hybrid representation: a budget of features picked from several base representations for one task and one model.
% Paragraph 2: the original motivation (one or two sentences)
% - The hybrid started as a fix for similarity: select the task-relevant features first, then compute distances only on those. It then turned out to work as an ordinary model input; this chapter returns to the original motivation.
% Paragraph 3: the methods under test (two or three sentences)
% - Nearest-neighbour editing \citep{Wilson1972}, filtering by classifiers trained on other folds \citep{Brodley1999}, and confident learning \citep{Northcutt2021}. Each depends on the feature space it works in.
% List: research questions, one per results section
% - Distances and nearest-neighbour prediction: is predicting each molecule's activity from its five nearest neighbours more accurate in a hybrid than in the single representations it is built from?
% - Noise detection and estimation: are noisy labels easier to find, or their size easier to estimate, in a hybrid?
% - Learning with noisy labels: does a hybrid change how well noise mitigation and noise-robust training work?
% Paragraph 4: status (one sentence)
% - Planned as the second hybrid paper. Methods are final; results are preliminary.
% - Result needed: one-sentence preview of the pattern found for each research question.
% ---- END NOTES ----


\section{Methods}
% ---- NOTES: Methods ----
% Takeaway: Every hybrid comes from the Chapter 5 search. Two things are new: a second score to rank candidate hybrids on, and what is done to the training labels.
% Takeaway: The hybrid is always rebuilt inside each fold from the training labels the experiment gives it, so feature choice sees the noise.
% Paragraph 1: what the experiments share
% - Datasets, scaffold-group-aware folds and metrics are as in Chapter \ref{chap:method}. Noise touches training labels only.
% - Comparison: each downstream model's own top two single representations, ranked by that model's own score.
% Paragraph 2: which labels the hybrid saw
% - The hybrid's features are rebuilt within each fold from that fold's training labels: noisy, clipped or uncleaned.
% - The search settings (representations, budget, allocation method) come from different labels per experiment:
%   - Survey of twelve LNL methods: searched again in each fold on the noisy training labels, with the same neural network the methods use. Budgets 50, 100 and 250 features.
%   - Labels clipped at a detection limit: settings from the search on the unclipped labels, which a real user would not have. State it as a limit.
%   - Cleaning real labels: settings from the search on the same real labels.
% Paragraph 3: the three tasks kept apart (two sentences)
% - Finding noisy labels, estimating how large the noise is, and improving the held-out score are separate tasks. A method can do the first well and fail the third.
% Figure: one-panel flow. Left: dataset with injected noise, or none. Middle: the hybrid built per fold on the training labels, beside the top two single representations. Right: similarity score, and the held-out score after each LNL method.
% ---- END NOTES ----


\subsection{Similarity in a hybrid}
% ---- NOTES: Similarity in a hybrid ----
% Takeaway: Two ways to use a hybrid for similarity: compute distances in an ordinary hybrid, or search a similarity hybrid whose candidates are ranked by nearest-neighbour prediction instead of a downstream model.
% Takeaway: In a hybrid, width is weight: a representation's pull on the distance is roughly its share of the budget.
% Paragraph 1: two ways to use a hybrid for similarity
% - Ordinary hybrid: the Chapter \ref{chap:method} search picks features that help a downstream model. Distances are then computed in that hybrid.
% - Similarity hybrid: the same search, but each candidate is ranked by how well each molecule's activity is predicted from its five nearest neighbours in the hybrid. No downstream model is fitted.
% - Everything else is shared: importance methods, allocation methods, budgets and folds. The neighbour score also sets each representation's share of the budget.
% - Every representation in the dataset's set is a candidate; the number of representations is not searched.
% Paragraph 2: the distance inside a hybrid
% - Manhattan distance (Chapter \ref{chap:flaws}) on features z-scored with the training molecules' mean and standard deviation.
% - Width is weight: every z-scored feature adds one term of similar size, so a representation's pull on the distance is roughly its number of features divided by the budget. The allocation method therefore also weights the representations inside the distance.
% - Result needed: confirm Manhattan distance and one-over-distance neighbour weighting once the design experiment on distance, neighbour weighting and number of neighbours has run.
% Paragraph 3: the neighbour score
% - Computed on the held-out molecules of each fold. Each molecule's 5 nearest neighbours come from other scaffold groups, because a same-scaffold neighbour is an easy match and inflates the score.
% - Prediction, weighting each neighbour by one over its distance: $\hat{y}_i = \frac{\sum_{j\in N_5(i)} y_j / d_{ij}}{\sum_{j\in N_5(i)} 1/d_{ij}}$.
% - Regression score: Spearman $\rho(\hat{\mathbf{y}}, \mathbf{y})$. Binary classification: $2(\mathrm{AUROC} - 0.5)$ on the distance-weighted fraction of neighbours in class 1, so chance is 0 under both.
% - The score measures how far the similar property principle (Chapter~\ref{chap:background}) holds in the hybrid.
% Paragraph 4: where the score is undefined (one or two sentences)
% - The similarity hybrid is not searched on three datasets: the PAM50 subtypes and trastuzumab HER2 multiclass (AUROC needs 2 classes), and trastuzumab HER2 binary (a held-out block can hold one edit-distance shell, so no molecule has a neighbour from another group).
% Figure: two pipelines sharing a first half (base representations -> importance -> allocation -> hybrid). Top: hybrid -> downstream model -> dataset metric. Bottom: hybrid -> 5 nearest neighbours from other scaffold groups -> distance-weighted mean activity -> Spearman or rescaled AUROC. Caption: same search, different score.
% ---- END NOTES ----
\todo[inline, color=yellow!40]{6.2.1: objective=similarity. [DIAGRAM D6.1]}

\paragraph{Feature importance weighting.}
% ---- NOTES: Feature importance weighting ----
% Takeaway: Scaling each feature by its importance before the distance lets features that matter for the task count more, and drops zero-importance features.
% Paragraph 1: what the weighted distance is (her paragraph below)
% - Change: her first sentence names noise detection as the only purpose. Add the second use: nearest-neighbour prediction, weighted against unweighted (results in "Results: similarity").
% - Importance is computed on each fold's training molecules only. Weights are applied after z-scoring; applied before, z-scoring divides them back out.
% - $w$ under Manhattan and $\sqrt{w}$ under Euclidean both multiply each feature's term in the distance by $w$ once.
% - Average 1 keeps weighted and unweighted distances on the same scale; neighbour rankings do not change under any rescaling of $w$.
% ---- END NOTES ----

Beyond simply using the feature importance scores for allocation, we investigate whether feature weighting by importance during distance computation improves noise detection. For a feature importance vector $w$, weighted distances are computed by scaling the standardised features: $x' = x \odot w$ under Manhattan distance ($x \odot \sqrt{w}$ under Euclidean), with negative importances set to 0 and $w$ rescaled to average 1. 


\subsection{Noise conditions}
% ---- NOTES: Noise conditions ----
% Takeaway: Three conditions drive the learning-with-noisy-labels experiments: censoring, Gaussian noise, and real labels with nothing injected. Outlier and grouped noise are used only to grade detection and estimation.
% Takeaway: Each condition stands for a real source of error, and noise is added to training labels only.
% Paragraph 1: which conditions this chapter uses
% - Defined in Section~\ref{sec:noise_injection} (Figure~\ref{fig:noise_conditions}). Noise is sized from each fold's clean training labels and added to training labels only.
% Paragraph 2: censoring at a detection limit
% - The level is the fraction of training labels clipped: 0 to 0.40. The clip has no random draw, so one run per level.
% - The error is one-sided, and its direction and location (the limit) are known.
% - Stands for assay range limits, e.g. "> 10 µM" recorded as 10 µM \citep{Svensson2025}. Regression only.
% Paragraph 3: Gaussian noise
% - Levels are multiples of the fold's clean training-label standard deviation, 0 to 1.5, three draws per level. On classification datasets, the fraction of labels flipped, 0 to 0.40.
% - Stands for the random error of repeating a measurement. It is also the standard assumption in noise-robustness studies.
% Paragraph 4: outlier, grouped-wider and grouped-shifted noise (detection and estimation only)
% - Outlier noise stands for annotation and transcription errors \citep{Kalliokoski2013}; here 1\%, 5\% and 10\% of training molecules.
% - Grouped-wider and grouped-shifted stand for differences between laboratories or sources \citep{Bentz2013}. Both need a scaffold per molecule, so run only on datasets with SMILES.
% Paragraph 5: real labels, nothing injected
% - No list of wrong labels exists, so success is only a better held-out score after cleaning. That shows cleaning helped the model, not that it found errors.
% - QM9 is the control: its labels are computed, with no experimental error \citep{Ramakrishnan2014}.
% Table 6.1: one row per condition. Columns: what it stands for; levels; which experiments use it; datasets.
% - Result needed: the datasets each condition covers in the full runs.
% ---- END NOTES ----
\todo[inline, color=yellow!40]{6.2.2: Table 6.1, one row per noise type.}


\subsection{Noise mitigation methods}
% ---- NOTES: Noise mitigation methods ----
% Takeaway: Three experiments, one per question: correcting labels clipped at a detection limit, a survey of twelve LNL methods under Gaussian noise, and cleaning real labels.
% Takeaway: Noise-robust training never has to say which direction a label is wrong in. Noise mitigation does, per molecule.
% Paragraph 1: the two kinds of LNL method
% - LNL, noise mitigation and noise-robust training are defined in Section~\ref{sec:lnl_background}. The first and third experiments are noise mitigation; the survey covers both.
% - Change: this heading says noise mitigation, but the survey also covers noise-robust training. A heading naming learning with noisy labels fits all three.
% Paragraph 2: correcting labels clipped at a detection limit
% - Training labels above a limit $c$ are recorded as $c$. The treatments read the limit off the training labels as the largest recorded value; they are never given it.
% - Reference prediction $p_i$: out-of-fold boosted-tree prediction, scaffold groups kept together. Residual spread $s$: SD of $y_i - p_i$ over molecules not at the limit.
% - Basis, the mean of a normal above a limit: $\mathbb{E}[Y \mid Y > c] = \mu + \sigma\frac{\phi(a)}{1-\Phi(a)}$, $a = (c-\mu)/\sigma$ \citep{Tobin1958}.
% - Treatments: per-molecule lift, $c + s\frac{\phi(a_i)}{1-\Phi(a_i)}$ with $a_i = (c-p_i)/s$; a shuffled control dealing those lifts to other clipped molecules; a constant lift, the mean of the per-molecule lifts; dropping the clipped molecules. Also run: the truncated-normal mean per molecule and winsorising.
% - Model comparators, not proposals (the thesis changes the input, not the model), each paired with the model it replaces: gradient boosting with a censored likelihood, a random forest with absolute-error splits, gradient boosting with the Huber loss.
% - Scored by recovery rate, defined below.
% Paragraph 3: a survey of twelve LNL methods
% - Question: does changing the input from a single representation to the hybrid change which methods work, and how well?
% - Base for every method: one plain fully connected network. Inputs: the top two single representations, the hybrid, and the importance-weighted hybrid.
% - Most of these methods were built for image classification with flipped labels. Here they are adapted to continuous labels.
% - Scored against the plain network on the same input and noise level.
% - Table: the twelve methods in three groups, one plain line each.
%   - Noise mitigation, remove then refit: small-loss selection \citep{Han2018}, confident learning \citep{Northcutt2021}, removal by distance to neighbours, staged removal \citep{zhao2025}.
%   - Noise-robust training that changes the training loop: per-epoch down-weighting (no source yet), curriculum by confidence \citep{Bengio2009}, DivideMix \citep{li2020dividemix}, Meta-Weight-Net \citep{shu2019metaweightnet}, MentorNet \citep{jiang2018mentornet}, Contrast to Divide \citep{zheltonozhskii2021contrast}.
%   - Noise-robust training that models the noise (continuous labels only): heteroscedastic GP \citep{liu2020}, deep evidential regression \citep{amini2020deep}.
% Paragraph 4: cleaning real labels, nothing injected
% - Per fold: fit on the uncleaned labels and score; clean the training labels; refit the same model and score again.
% - Two suspicion scores: neighbour residual (distance from the mean label of 10 nearest neighbours in other scaffold groups) and cross-validation residual (distance from an out-of-fold prediction).
% - Seven cleanups, one line each: remove flagged molecules; down-weight by suspicion; move every label toward its neighbours' mean in proportion to suspicion; move only flagged labels there; move labels toward an out-of-fold prediction; toward a Gaussian-process mean; clip labels to the range carried by similar training molecules.
% - Each is run at a few strengths (share flagged, or how far labels move).
% - Shuffled-score control: the same cleanups with suspicion scores dealt at random. It separates "which molecules" from "how many".
% - Scored as folds where the cleaned model beats the uncleaned one, out of folds run.
% Paragraph 5: what the four subsubsections below are for (one sentence)
% - The three experiments use only the neighbour residual, the cross-validation residual, proportional smoothing, inverse weighting and the recovery rate. The rest describes the exploratory study (Section~\ref{sec:lnl_motivation}).
% Figure: true label against recorded label for the clipped-label experiment, with the pile-up at $c$ and arrows for the lift treatments.
% ---- END NOTES ----
\todo[inline, color=yellow!40]{6.2.3: methods for censoring correction, the LNL survey and cleaning real labels.}

In addition to evaluating hybrid representation quality for distance-based reasoning, our investigation evolved to explore alternative noise handling strategies. Initial experiments with sample removal methods revealed limitations when applied to Gaussian noise affecting all samples. This motivated a pivot toward sample weighting and label smoothing approaches, which do not require accurate binary classification of noisy versus clean samples but instead leverage continuous noise magnitude estimates.

\subsubsection{Sample removal methods}
% ---- NOTES: Sample removal methods ----
% Takeaway: Removal was the first approach tried in the exploratory study. Under Gaussian noise every label is wrong by some amount, so removal throws away labels that are only partly wrong. This led to weighting and smoothing.
% Paragraph 1: the exploratory study (two sentences)
% - QM9 HOMO-LUMO gap, 10,000 molecules, one scaffold split; Gaussian noise of SD 0.3 and 0.6 eV on training labels; ECFP4, PDV and MHG-GNN; random forest after removal.
% - Cite only the single-representation rows. Some hybrid rows used an allocation method no longer in the thesis.
% Paragraph 2: k-nearest-neighbour filtering
% - Change: the first sentence re-states the similar property principle with a different citation from Chapter~\ref{chap:background}. Cite \citet{Johnson1990} as that chapter does.
% - Change: the sentence beginning "Presumably" can state the idea directly. Editing a training set by its nearest neighbours goes back to \citep{Wilson1972}.
% Paragraph 3: Local Outlier Factor
% - Change: there is no "LOF greater than 2" rule. LOF uses 20 neighbours and flags the 10\% of training molecules with the highest score.
% - LOF sees only the features, never the labels, so it flags molecules in sparse regions, not molecules with unusual labels.
% Paragraph 4: the other removal methods
% - Change: rename "AC detection" to "activity-cliff detection".
% - Co-teaching \citep{Han2018} and DivideMix \citep{li2020} were cut down to their selection step, run with random forests, and their rejected molecules deleted. Say so, so a reader does not expect the original methods.
% - Cross-validation filtering follows \citep{Brodley1999}.
% Paragraph 5: why removal does not suit Gaussian noise (one or two sentences)
% - Under Gaussian noise every training label is wrong by some amount, so detection precision is 1 by construction. The useful question becomes how wrong each label is.
% ---- END NOTES ----

Distance-based approaches rely on the assumption that similar molecules have similar properties \citep{Cherkasov2014}. Presumably, noisy data would break this pattern, making it a method for detecting anomalous data. k-Nearest neighbor filtering examines the variance of a sample's neighborhood. For sample $x_i$ with $k$ nearest neighbors $\mathcal{N}_k(x_i)$, we compute the standard deviation of neighbor labels:
$$
\sigma_i = \text{std}(\{y_j : x_j \in \mathcal{N}_k(x_i)\})
$$
Samples with $\sigma_i$ exceeding the 90th percentile are marked as potentially noisy. We test $k \in \{5, 10, 20\}$.

Local Outlier Factor (LOF) identifies anomalous samples by detecting local density estimates between a sample and its neighbors \citep{Breunig2000}. These density estimates are based on the distance to its $k$ nearest neighbors. If a point's density differs from that of its neighbors such that the LOF is substantially greater than 2, it is flagged as an outlier.

We also tested the following noise removal methods:  distance-weighted prediction filtering, neighborhood consensus filtering, AC detection, Mahalanobis outlier detection, cross-validation disagreement, bootstrap ensemble filtering, co-teaching \citep{Han2018}, and DivideMix \citep{li2020}.

\subsubsection{Noise magnitude estimation}
% ---- NOTES: Noise magnitude estimation ----
% Takeaway: A score of how wrong each label is lets later steps down-weight or smooth in proportion. Removal is all or nothing.
% Paragraph 1: why score how noisy each label is
% - \citet{Martin2021} adapt classification noise filters to continuous labels and find they help regression models. Cite it for the general idea only.
% - Change: the first sentence calls estimation "the core of the noise mitigation research". Estimating noise is not itself noise mitigation; it is the step noise mitigation builds on.
% Paragraph 2: scores from each molecule's nearest neighbours
% - Neighbour residual is the continuous-label form of editing a training set by neighbour disagreement \citep{Wilson1972}.
% - Change: the prose writes the local variance ratio's denominator as std + epsilon. It is a floor: the neighbours' SD, never below 1\% of the dataset's label SD.
% - Change: spell out "activity cliff" instead of "AC" \citep{Stumpfe2012, Dablander2023}.
% Paragraph 3: scores from models that never saw the molecule
% - Cross-validation residual: cite \citep{Brodley1999}; scaffold groups kept together.
% Paragraph 4: scores from models that predict their own uncertainty
% - All are fitted out of fold, so no molecule is scored by a model that trained on it.
% - Change, Gaussian process: one learned noise level is shared by every molecule, so it cannot rank molecules. The score is the out-of-fold predictive SD plus the out-of-fold absolute residual.
% - Change, random forest uncertainty: the prose calls it "Bayesian random forest" and says out-of-bag. It is the spread across trees plus half the out-of-fold absolute residual \citep{taddy2015}.
% - Change: the last sentence says NGBoost was "selected". It is one score among the others. The correlation it cites from Section~\ref{sec:uq} is pooled across noise levels; within one noise level it is near zero.
% - Five more scores are computed but not in the prose: distance-weighted neighbour residual, out-of-bag residual, combined score, heteroscedastic GP, quantile regression forest.
% Table (optional): one row per noise score; columns family, what is computed in one line, citation. It can replace the three itemised lists.
% ---- END NOTES ----

We now examine the core of the noise mitigation research: estimating the magnitude of noise. This enables proportional handling rather than complete removal. We used the following neighbor-based estimators:

\begin{itemize}
    \item Neighbor residual: Estimates noise as deviation from k-NN mean, normalized by global label standard deviation:
    $$\text{score}_i = \frac{|y_i - \bar{y}_{\mathcal{N}_k(x_i)}|}{\sigma_y}$$
    
    \item Local variance ratio: Measures deviation relative to local label variance, identifying samples that deviate from neighbors more than neighbors deviate from each other:
    $$\text{score}_i = \frac{|y_i - \bar{y}_{\mathcal{N}_k(x_i)}|}{\text{std}(y_{\mathcal{N}_k(x_i)}) + \epsilon}$$
    
    \item AC-aware: Distinguishes noise from activity cliffs by weighting deviation by neighborhood consistency. High deviation with high neighbor agreement suggests noise while high deviation with neighbor disagreement suggests an AC \citep{Dablander2023}:
    $$\text{score}_i = \frac{|y_i - \bar{y}_{\mathcal{N}_k(x_i)}|}{\sigma_y} \cdot \frac{1}{1 + \text{std}(y_{\mathcal{N}_k(x_i)})/\sigma_y}$$
\end{itemize}

And we used the following model-based estimators using prediction disagreement:

\begin{itemize}
    \item Cross-validation residual: Uses out-of-fold predictions; high residuals indicate labels inconsistent with the model learned from other samples.
    \item Ensemble disagreement: Trains multiple models and combines prediction variance with mean residual.
\end{itemize}

We used probabilistic methods that directly predict noise magnitude using uncertainty estimates:

\begin{itemize}
    \item Gaussian Process: Provides natural uncertainty estimates via posterior variance, with learned noise level $\sigma^2$ from the kernel \citep{gauche}.
    
    \item Heteroscedastic Neural Network: Predicts both $\mu(x)$ and $\sigma^2(x)$ using $\beta$-NLL loss \citep{Seitzer2022}:
    $$\mathcal{L} = \frac{1}{2}\left(\frac{(y - \mu)^2}{\sigma^2} + \log \sigma^2\right) \cdot \sigma^{2\beta}$$
    where $\beta = 0.5$ prevents the variance from hiding the mean errors.
    
    \item Conformal prediction: Uses split conformal prediction to obtain prediction intervals. Wider intervals should indicate higher uncertainty \citep{Shafer2008}.
    
    \item Deep ensemble: Trains $M$ independent networks and uses the variance of means as epistemic uncertainty plus the mean of variances as aleatoric uncertainty \citep{lakshminarayanan2017}.
    
    \item Bayesian Random Forest: Uses tree prediction variance and out-of-bag residuals as uncertainty estimates \citep{taddy2015}.
    
    \item Quantile regression: Fits models for multiple quantiles and uses interval width as noise estimate.
    
    \item NGBoost: Natural gradient boosting that predicts distributional parameters $(\mu, \sigma)$ directly \citep{Duan2020}.
\end{itemize}

NGBoost in particular was selected due to its promising uncertainty values with respect to the correlation between noise and prediction error, as discussed in Section~\ref{sec:uq}.  

\subsubsection{Label smoothing and sample weighting}
% ---- NOTES: Label smoothing and sample weighting ----
% Takeaway: Smoothing moves a label part of the way toward a target; weighting keeps the label and lowers how much the molecule counts. Only the proportional rule and the inverse weight are used in this chapter's experiments.
% Paragraph 1: before the lists, define the two inputs once
% - score_norm: the noise score rescaled to 0-1 over the fold's training molecules. Target: the mean label of the nearest training neighbours, or a model's prediction. alpha: strength, 0 (no change) to 1 (label replaced).
% Paragraph 2: changes to the prose
% - Threshold rule: the top alpha share by score is flagged, and each moves alpha of the way to its target, so alpha sets both how many labels move and how far.
% - Soft threshold: add the second line, $y + \text{strength}\cdot(\text{target} - y)$.
% - "to preserve effective sample size" is wrong. Say "so the total weight equals the number of training molecules".
% ---- END NOTES ----

Given the fact that all labels contain some magnitude of noise, we focused strongly on label smoothing over removal using the following methods:

\begin{itemize}
    \item Proportional: $y_{\text{smooth}} = y + \alpha \cdot \text{score}_{\text{norm}} \cdot (\text{target} - y)$
    
    \item Threshold: Only correct samples above the $(1-\alpha)$ percentile of scores.
    
    \item Soft threshold: Applies sigmoid-like correction strength:
    $$\text{strength} = \alpha \cdot \sigma\left(5 \cdot (\text{score}_{\text{norm}} - 0.5)\right)$$
    
    \item Confidence: For methods that output raw uncertainty $\sigma$ (e.g., NGBoost):
    $$\text{correction} = \frac{\alpha \cdot \sigma}{1 + \alpha \cdot \sigma}, \quad y_{\text{smooth}} = (1 - \text{correction}) \cdot y + \text{correction} \cdot \mu$$
\end{itemize}


Rather than directly modifying labels, we weighted more trustworthy samples during training:

\begin{itemize}
    \item Inverse: $w = \frac{1}{1 + \alpha \cdot \text{score}_{\text{norm}}}$
    
    \item Exponential: $w = \exp(-\alpha \cdot \text{score}_{\text{norm}})$
    
    \item Soft: $w = 1 - \alpha \cdot \text{score}_{\text{norm}}^2$, clipped to $[0.1, 1.0]$
\end{itemize}

All weights are normalized to mean 1 to preserve effective sample size.

\subsubsection{Recovery rate}
% ---- NOTES: Recovery rate ----
% Takeaway: Recovery is the share of the score lost to label noise that an LNL method gives back. 100\% restores clean performance; 0\% is no change; negative is worse than leaving the noise in.
% Takeaway: Detection is scored separately, with three numbers that need no fixed cut-off.
% Paragraph 1: how it is computed here
% - Change: the prose describes an older protocol, one random forest on ECFP4 scored by $R^2$. Replace the first-person opening ("We have a standard RF...") with a plain description: the same model fitted on clean, noisy and treated training labels, test labels never altered.
% - Every downstream model is scored in the dataset's own metric and divided by its own clean-to-noisy gap. The formula's direction cancels, so it works for RMSE too.
% - No recovery is reported when the gap is under 2\% of the clean score, or when the noisy score is no worse than the clean one.
% - Change: use \[ ... \] for the display equation instead of $$ ... $$.
% Paragraph 2: how detection is scored
% - Change: the prose says precision, recall and F1. The experiments report average precision (the headline), precision in the top k (k = number of changed molecules), and AUROC.
% - Detection is not reported under a noise type that changes every label.
% ---- END NOTES ----

For noise detection accuracy, we compute precision, recall, and F1 score, treating noise injection as ground truth.

We have a standard RF on ECFP4 train and test on the clean dataset, producing $R^2_{\text{clean}}$. Noise is then added to the clean data, RF on ECFP4 is then freshly trained and tested on this noisy data, producing $R^2_{\text{noisy}}$. Using the given estimation method and representation to test, noise estimation and label smoothing is applied to the noisy data, resulting in the denoised set. RF on ECFP4 is yet again run, this time on the denoised data, to produce $R^2_{\text{denoised}}$. From these performance values, we derive recovery rate, or the fraction of performance loss recovered after denoising, as follows: 

$$
\text{Recovery} = \frac{R^2_{\text{denoised}} - R^2_{\text{noisy}}}{R^2_{\text{clean}} - R^2_{\text{noisy}}}
$$

A recovery rate of 1.0 indicates complete restoration to clean performance; negative values indicate the denoising procedure made performance worse than the noisy baseline.


\section{Results: similarity}
% ---- NOTES: Results: similarity ----
% Takeaway: Answers the first research question: is predicting each molecule's activity from its nearest neighbours more accurate in a hybrid, or in a similarity hybrid, than in the single representations?
% Takeaway: Also asks whether scaling features by importance gives neighbours with closer activities.
% Paragraph 1: nearest-neighbour prediction in a hybrid against the single representations
% - Prediction: the mean activity of the 5 or 10 nearest training molecules.
% - Spaces compared: each single representation, the top three ordinary hybrids and the top three similarity hybrids, each hybrid with and without importance weighting.
% - Reference: for each held-out fold, the best single representation and best hybrid are chosen on the other folds. Report hybrid minus that representation, in the dataset's own metric, per fold.
% - Result needed: held-out score per dataset and fold for the best single representation, the ordinary hybrid and the similarity hybrid.
% Table: one row per dataset. Columns: best single representation (name and score), ordinary hybrid, similarity hybrid, each minus the single representation. Entry: mean over folds with the fold range.
% Paragraph 2: distance against activity difference
% - For sampled pairs of held-out molecules, absolute activity difference against distance, laid out like the hERG Ki figure in Chapter \ref{chap:flaws}.
% - Two panels: best single representation and similarity hybrid. One number per panel: Spearman correlation between distance and activity difference.
% - Result needed: the figure, and the correlation per space for each dataset.
% Figure: two scatter panels side by side, activity difference (dataset units) against distance. Caption gives the correlation for each panel.
% Paragraph 3: importance-weighted distance for k-nearest-neighbour prediction
% - Space: all of the dataset's representations joined side by side, not a hybrid. Each feature multiplied by the square root of its importance; Euclidean distance; k = 5, 10 and 20.
% - Scored in $R^2$ on regression datasets, not the dataset's own metric. One sentence.
% - No control shuffles the weights across features, so a gain cannot be split into "which features" and "unequal weights".
% - Result needed: unweighted and weighted $R^2$ per dataset, importance method and k.
% Paragraph 4: the similarity hybrid as an ordinary model input
% - In the cleaning of real labels, downstream models are also fitted on the similarity hybrid. With no cleaning, this compares the two hybrids as inputs to the same model on the same folds.
% - Result needed: held-out score with no cleaning, per dataset, model and fold, for the similarity hybrid, the ordinary hybrid and the top two single representations.
% ---- END NOTES ----
\todo[inline, color=yellow!40]{Results TODO(prelim): methods are final, numbers will be replaced by the server runs. [DIAGRAM D6.2]}


\section{Results: noise detection and estimation}
% ---- NOTES: Results: noise detection and estimation ----
% Takeaway: Two of the three tasks: finding which training labels are wrong (detection), and sizing how wrong they are (estimation).
% Takeaway: Finding noise is possible for outliers, and a systematic offset can be sized; correcting labels to improve the held-out score is a different problem.
% Paragraph 1: how detection and estimation are scored
% - Each noise score ranks a fold's training molecules from most to least suspect. Detection is graded against the molecules the injection singled out: widened under outlier and grouped-wider noise, clipped under censoring.
% - Chance level is the share singled out, e.g. about 0.1 at 10\% outlier noise. Report it beside every average precision.
% - Estimation is graded against the injected noise size in label units, or a size measured from repeat measurements.
% Paragraph 2: finding injected noisy labels, and whether the space matters
% - Question: does a neighbour-based score find the corrupted molecules more often when neighbours are found in a hybrid than in a single representation?
% - Ways to find neighbours: the single representation, the same weighted by importance, forest proximity, the hybrid, and the similarity hybrid. Model-based scores change only their input features.
% - Result needed: average precision per score, space and noise condition, with chance level and spread across folds.
% - Preliminary, will be replaced: antibody thermostability, outlier noise on 10\% of labels, one-hot, no hybrid. Average precision: plain Euclidean 0.1905, TreeSHAP-weighted Euclidean 0.2292, forest proximity 0.2251; chance about 0.09 (about 50 moved of 537 training molecules).
% - Preliminary, will be replaced, unverified: hERG Ki, outlier noise on 10\% of labels, one fold. Ten single representations reach average precision 0.4036 to 0.4326; every way of combining the ten falls inside that range.
% Paragraph 3: finding real suspect labels, nothing injected
% - Curator flags: ChEMBL curators mark suspect records \citep{Zdrazil2023}. On hERG IC50 in ChEMBL 37, 301 of 12,008 records carry a flag, on 215 of 9,802 molecules.
% - Most are flagged "outside typical range" because the value is extreme, so a score reading distance from neighbours is partly graded on the value itself. State this next to the number.
% - Inter-laboratory disagreement: BindingDB holds the same compound measured by different laboratories \citep{Gilson2016}. The Spearman correlation between a score and that spread grades it against real noise.
% - Result needed: average precision against curator flags; correlation with inter-laboratory spread.
% Paragraph 4: estimating how large the noise is
% - From repeat measurements: a molecule's label error variance is the per-measurement variance divided by its number of measurements $k$. Regressing squared out-of-fold residuals on $1/k$ gives the per-measurement noise variance as the slope.
% - From near-duplicate pairs: half the mean squared label difference of the closest 1\% of training pairs. An upper bound.
% - From the gamma test \citep{Stefansson1997}. Also an upper bound.
% - Put each estimate beside the model's own held-out mean squared error. Noise mitigation can only pay where label error is a large share of model error. Compare with the published size of Ki measurement error \citep{Kramer2012}.
% - Result needed: the three estimates per dataset and space, beside the model's mean squared error.
% Paragraph 5: the gap between finding noise and correcting it
% - Flag sets are built with precision set by hand, 0 to 1 in steps of 0.25, and each is acted on. This gives the precision a score would need before acting improves the held-out score. Compare it with the average precision the scores reach.
% - Result needed: the precision needed before acting pays, per dataset, against the precision reached.
% Table: one row per noise score. Columns: average precision under outlier noise at 10\%, grouped-wider noise and censoring, each in the best single representation and in the hybrid; chance level in the header row.
% ---- END NOTES ----
\todo[inline, color=yellow!40]{Results TODO(prelim): methods are final, numbers will be replaced by the server runs.}

\section{Results: noise mitigation by noise type}
% ---- NOTES: Results: noise mitigation by noise type ----
% Takeaway: Results go experiment by experiment (clipped labels, the twelve-method survey, cleaning real labels), then one closing subsection compares what each noise type allows.
% Paragraph 1: what every result is compared against
% - Each model's own top two single representations, the hybrid, and, where run, the importance-weighted or similarity hybrid. Scores in each dataset's own metric, per fold, no significance tests.
% - Change: the section heading says noise mitigation, but the survey also covers noise-robust training.
% - Caveat for the preliminary numbers: they compare against one fixed single representation per dataset, one-hot on antibody thermostability (its weakest representation for a random forest, 0.3471 Spearman) and ECFP4 on hERG Ki and Caco-2.
% ---- END NOTES ----


\subsection{Censoring: correcting labels clipped at a detection limit}
% ---- NOTES: Censoring: correcting labels clipped at a detection limit ----
% Takeaway: Raising every clipped label by one constant recovered most of the damage from clipping; per-molecule amounts did no better, and deleting the clipped molecules was worse than keeping them (preliminary).
% Paragraph 1: the question and why the clip is injected
% - Question: when training labels past an assay's limit are recorded at the limit, can they be corrected so the model predicts held-out molecules better?
% - Injected, not real: the hERG Ki data keep only ChEMBL records with an exact "=" relation, so clean training labels are clipped on purpose. Held-out labels keep their real values.
% Paragraph 2: hERG Ki at 10\% clipped (preliminary, will be replaced)
% - Setup: ECFP4, 3 folds, random forest, LightGBM, kNN and SVM. Limit 6.86 pKi, about 90 of 941 training molecules per fold.
% - Deleting the clipped molecules scores worse than keeping them clipped, on all four models, in every fold and at every clipping level.
% - Robust losses barely move from the clipped row: random forest with absolute-error splits 0.6574 pKi, LightGBM with the Huber loss 0.6480 pKi.
% Paragraph 3: other clipping levels and the per-molecule lift (preliminary, will be replaced)
% - Also run: hERG Ki at 25\% and 40\% clipped, and Caco-2 efflux ratio at 25\%.
% - The constant lift beat the clipped labels in 48 of 48 comparisons (4 models, 3 folds, 4 conditions). Recovery shrinks as more labels are clipped.
% - A separate lift per molecule scored higher RMSE than the constant in 15 of 16 model-condition pairs. Dealing the per-molecule lifts to the wrong molecules scored lower RMSE than dealing them correctly, in 16 of 16.
% - So per-molecule lift sizes did no better than one shared size, and no better than the same sizes given to the wrong molecules.
% - Result needed: recovery rate per treatment, model, feature space and clipping level; the three model comparators against their parent models.
% Table: hERG Ki, 10\% of training labels clipped at 6.86 pKi, ECFP4, RMSE (pKi), mean of 3 folds. One row per version of the training labels; columns random forest, LightGBM, kNN, SVM.
% - Never clipped: 0.6032, 0.6087, 0.6871, 0.6841.
% - Clipped: 0.6588, 0.6499, 0.7158, 0.7148.
% - Clipped, then constant lift: 0.6118, 0.6248, 0.6937, 0.6847.
% - Clipped molecules deleted: 0.7706, 0.7631, 0.8037, 0.8162.
% Table: recovery rate of the constant lift, ECFP4, mean of 3 folds. One row per dataset and clipping level; columns random forest, LightGBM, kNN, SVM.
% - hERG Ki 10\%: 85\%, 61\%, 77\%, 98\%.
% - hERG Ki 25\%: 56\%, 62\%, 54\%, 65\%.
% - hERG Ki 40\%: 45\%, 48\%, 51\%, 53\%.
% - Caco-2 25\%: 65\%, 71\%, 72\%, 68\%.
% - Print the raw RMSE pairs beside each share, or in the Additional files.
% ---- END NOTES ----
\todo[inline, color=yellow!40]{Results TODO(prelim): methods are final, numbers will be replaced by the server runs.}


\subsection{Gaussian noise: learning-with-noisy-labels methods}
% ---- NOTES: Gaussian noise: learning-with-noisy-labels methods ----
% Takeaway: Under Gaussian noise the input moved the score more than the method: the hybrid with no LNL method beat the best method on one-hot (preliminary).
% Takeaway: The methods helped only at the noise level that actually lowered the plain network's score.
% Paragraph 1: what was run (preliminary, will be replaced)
% - Antibody thermostability, 672 antibodies, melting temperature. Spearman on test antibodies, mean of 3 folds.
% - Gaussian noise at 0.5 and 1.5 times the SD of the fold's clean training labels.
% - Inputs: one-hot (5600 features), the hybrid built on each fold's noisy labels, and the importance-weighted hybrid.
% Paragraph 2: the plain network before any method
% - No added noise: one-hot 0.4167, hybrid 0.5418, importance-weighted hybrid 0.5857 Spearman.
% - Noise at 0.5 times did not lower one-hot or the hybrid. At 1.5 times the hybrid fell from 0.5418 to 0.4859. So 1.5 times is the level that left damage to repair.
% Paragraph 3: how many of the twelve methods beat the plain network on the same input
% - At 0.5 times: 0 of 12 on one-hot, 0 of 12 on the hybrid, 5 of 12 on the importance-weighted hybrid.
% - At 1.5 times: 8 of 12 on one-hot, 8 of 12 on the hybrid, 6 of 12 on the importance-weighted hybrid.
% Paragraph 4: input against method
% - The hybrid scored above one-hot at both levels, for the plain network and every method.
% - At 1.5 times, the best method on one-hot (small-loss selection, 0.3568) is below the plain network on the hybrid (0.4859). The two ranges do not overlap.
% Paragraph 5: limits (two sentences)
% - 3 folds, one noise draw per level, and the hybrid was searched with a random forest while the methods train a network. The final runs fix both and replace one-hot with each model's top two single representations.
% - Result needed: per noise level and dataset, the count of methods beating the plain network on each input, and whether the input moves the score more than the method.
% Table: one row per method, the plain network first (0.4534, 0.5991, 0.5833 at 0.5 times; 0.2729, 0.4859, 0.5118 at 1.5 times), then the twelve in the order of the methods table. Columns: one-hot, hybrid and importance-weighted hybrid at 0.5 times, then the same three at 1.5 times. Cells: Spearman, mean of 3 folds. A line above the table says what one row is.
% ---- END NOTES ----
\todo[inline, color=yellow!40]{Results TODO(prelim): methods are final, numbers will be replaced by the server runs.}


\subsection{Real labels, nothing injected: cleaning}
% ---- NOTES: Real labels, nothing injected: cleaning ----
% Takeaway: On the two small-molecule datasets no cleanup improved on leaving the labels as they came (preliminary).
% Takeaway: On the antibody dataset whether cleaning helped depended on the feature space, not only on the cleanup (preliminary).
% Paragraph 1: what was run (preliminary, will be replaced)
% - hERG Ki and Caco-2 (RMSE, lower is better) and antibody thermostability (Spearman). Random forest, 3 folds; each number averaged over every setting of that cleanup.
% Paragraph 2: the two small-molecule datasets
% - hERG Ki, RMSE in pKi. No cleaning: ECFP4 0.6032, hybrid 0.6029, similarity hybrid 0.6235. Every cleanup, in every space, scored at or above no cleaning.
% - Caco-2, RMSE in log efflux ratio. No cleaning: ECFP4 0.3242, hybrid 0.3090, similarity hybrid 0.3226. Down-weighting matched no cleaning; every other cleanup was worse.
% Paragraph 3: the antibody dataset, where the direction depends on the feature space
% - Spearman. No cleaning: one-hot 0.4068, hybrid 0.6279, similarity hybrid 0.6087.
% - On one-hot every cleanup lowered the score (worst: removal, 0.3498). On the hybrid every cleanup raised it (largest: removal 0.6429). On the similarity hybrid the seven split three above, four below.
% - Do not call it a gain until the shuffled-score control is in. The range clip reads no suspicion score, so that control cannot test it.
% - Result needed: the shuffled-score control per cleanup and feature space, antibody dataset first.
% Paragraph 4: what the final runs report
% - Folds where the cleaned model beat the uncleaned one, out of folds run, with the mean difference beside it. QM9 reported apart as the control.
% - Result needed: folds won per cleanup, suspicion score, feature space and dataset, including QM9.
% Table: one block per dataset. Rows: no cleaning, then the seven cleanups. Columns: top two single representations, the hybrid, the similarity hybrid. Cell: folds won out of folds run, with the mean difference in the dataset's metric. The line above the table says what one row is and what one cell counts.
% ---- END NOTES ----
\todo[inline, color=yellow!40]{Results TODO(prelim): methods are final, numbers will be replaced by the server runs.}


\subsection{By noise type}
% ---- NOTES: By noise type ----
% Takeaway: Gaussian noise is the standard assumption and the hardest to fix: every label is wrong by an unknown amount in an unknown direction, and the neighbours used to judge it are noisy too.
% Takeaway: Noise can be found or sized and still not be removable. Better similarity does not solve this.
% Takeaway: Ask what kind of noise is likely before choosing a method.
% Paragraph 1: why Gaussian noise is the hardest to fix
% - Discussion, no new results.
% - Every label carries an error of unknown size and sign. No clean subset exists to learn from.
% - A detector that compares a molecule with its neighbours compares it with labels that are also noisy. A large disagreement can be noise or a real activity cliff.
% - Even a correctly flagged label has no known correction. Removing it loses a molecule. Replacing it with the neighbours' mean swaps its noise for theirs.
% Paragraph 2: finding the noise is not the same as removing it
% - Noise can be identified or estimated, yet a cleaned dataset can give a worse model than the noisy one it came from.
% - Evidence: the exploratory study ran 860 cleaning configurations on the QM9 HOMO-LUMO gap (Chapter~\ref{chap:flaws}). Only 8 beat training on the noisy labels at both noise levels (SD 0.3 and 0.6 eV).
% - The survey's preliminary numbers do not yet support this: at 0.5 times the noise did no damage to repair, and at 1.5 times 8 of 12 methods helped.
% - A better similarity can make neighbours closer in true activity. It cannot remove the noise on their labels.
% Paragraph 3: the other noise types, from easiest to fix to hardest
% - Censoring is the most fixable: the limit and the direction are known, only the size is not. That is why one constant lift recovered most of the damage.
% - Outlier noise is easier to find: the widened labels stick out against their neighbours. The fix is removal or down-weighting; neither recovers the true value.
% - Grouped-wider noise makes every neighbour in the family noisy too. Grouped-shifted noise moves a family together, so comparing with same-family neighbours cannot reveal it; knowing each record's source makes the offset estimable.
% - These three have detection and estimation results but no LNL result in this chapter. One sentence.
% Paragraph 4: practical takeaways
% - Data pooled from many laboratories: expect an offset per source. Keep the source of every record.
% - A single controlled experiment: expect censoring. Keep the "<" and ">" qualifiers.
% - Molecules from very different chemical domains: neighbour-based detection is least reliable.
% - Source and amount unknown: noise-robust training is safer than noise mitigation (Chapter~\ref{chap:discussion}).
% - Result needed: per noise type, whether any LNL method improved on the untreated model, on each feature space.
% Table: one row per noise type (Gaussian, outlier, grouped-wider, grouped-shifted, censoring). Columns: what is wrong in the labels; how easy the noisy labels are to find; whether the direction of the error is known; what makes a correction possible.
% Figure: five panels, one per noise type. x is the true label, y the observed label; the diagonal is no noise. Shade the part that can be detected; mark the part that can be corrected.
% ---- END NOTES ----
\todo[inline, color=yellow!40]{6.5.5: why Gaussian is hardest; practical takeaways. [DIAGRAM D6.3]}


\section{Conclusions}
% ---- NOTES: Conclusions ----
% Takeaway: One paragraph per research question, then what each noise type allows, the limits, and the bridge: model the noise before choosing a method.
% Takeaway: Short, 1-2 pages, no formulas.
% Paragraph 1: distances and nearest-neighbour prediction
% - Result needed: per dataset, nearest-neighbour prediction in the ordinary and similarity hybrid minus the best single representation; and the similarity hybrid as a downstream-model input.
% - If nearest-neighbour prediction improves when features are chosen for the property, then "similar" depends on the property, not on the two structures alone. Write this only if the results show it.
% Paragraph 2: noise detection and estimation
% - Finding a noisy label and correcting it are separate problems.
% - Result needed: per dataset, detection and estimation in the hybrid against the top two single representations.
% Paragraph 3: learning with noisy labels (one sentence per experiment)
% - Result needed: censoring, Gaussian-noise survey, and real-label cleaning, each against leaving the labels alone.
% Paragraph 4: what each noise type allows (two sentences)
% - The type of noise decides whether it can be corrected, not only found. Censoring has a known location and direction; Gaussian noise has neither.
% Paragraph 5: limit specific to this chapter
% - The clipped-label experiment takes its search settings from the unclipped labels, which a real user would not have.
% Paragraph 6: bridge to Chapter \ref{chap:discussion}
% - Model the noise before choosing a method. Which method helps depends on noise type, amount, dataset and feature space.
% ---- END NOTES ----


